%*******************************************************
% Introduction
%*******************************************************

\chapter*{Introduction}\label{ch:introduction}
\addcontentsline{toc}{chapter}{Introduction}

\section*{La transition numérique comme question énergétique}

Depuis deux décennies, la transition numérique est présentée comme un vecteur de décarbonation et de croissance dématérialisée. Les technologies de l'information et de la communication optimiseraient les processus industriels, réduiraient les besoins de déplacement, rendraient les réseaux électriques intelligents. Cette promesse s'accompagne d'un récit implicite~: le numérique serait essentiellement immatériel, et son empreinte énergétique, négligeable au regard des gains d'efficacité qu'il procure.

Or les faits résistent à ce récit. La consommation électrique des datacenters croît désormais 2,5 fois plus vite que la demande totale d'énergie. Les modèles d'intelligence artificielle à grande échelle requièrent des puissances de calcul qui doublent tous les quelques mois. Les géants du numérique signent des contrats d'approvisionnement en énergie nucléaire pour alimenter leurs infrastructures. Le numérique n'est pas immatériel~: il est une industrie lourde de l'information, dont le substrat est l'électricité.

Cette thèse part de ce constat pour poser une question à la fois simple et structurante~: quelles sont les dimensions économiques et énergétiques de l'infrastructure ICT, et comment les représenter dans les outils d'analyse et de modélisation~?

\section*{Les enjeux de représentation}

La réponse à cette question se heurte à un problème de représentation. Les modèles économiques traitent le plus souvent l'ICT comme un facteur de productivité exogène~: un chiffre qui entre dans une fonction de production, un paramètre qu'on calibre sur des données passées. Les modèles énergie-climat, de leur côté, représentent la consommation des datacenters comme un poste de demande électrique parmi d'autres, sans endogénéiser les mécanismes qui font croître cette demande.

Entre les deux, un angle mort~: personne ne modélise le couplage structurel entre le calcul et l'énergie. Personne ne représente le fait que la loi de Moore, en améliorant l'efficacité unitaire du calcul, a paradoxalement multiplié la demande agrégée d'énergie~--- un paradoxe de Jevons computationnel que les cadres analytiques existants ne capturent pas.

Cette thèse propose de combler cet angle mort. Pour cela, elle établit d'abord l'existence et la nature du couplage ICT-énergie en le regardant depuis trois disciplines~--- l'histoire des techniques, la macroéconomie, l'économie de l'énergie~--- puis elle le modélise dans un cadre énergie-économie-climat.

\section*{Questions de recherche}

La question générale se décline en plusieurs sous-questions~:

\begin{itemize}
    \item Comment le couplage entre ICT et énergie s'est-il construit historiquement, et pourquoi est-il devenu structurel après 1971~?
    \item Quels sont les canaux économiques par lesquels l'ICT transforme la structure productive, et quelles sont les implications pour la demande d'énergie~?
    \item Les promesses de décarbonation par l'ICT résistent-elles à l'examen empirique~? Quelles sont les barrières systémiques~?
    \item Comment endogénéiser le calcul dans un modèle énergie-économie-climat tel qu'IMACLIM~?
    \item Quels scénarios émergent lorsqu'on distingue un déploiement ICT volontariste et coordonné d'un déploiement au fil de l'eau~?
\end{itemize}


\section*{Analyser le changement technique de grande ampleur}

Mais avant d'examiner ces cadres, il faut poser un fait préalable~: l'objet même qu'ils prétendent analyser~--- «~l'ICT~»~--- n'existe comme catégorie normée que depuis 1998, date à laquelle les pays de l'OCDE s'accordent sur une définition sectorielle des «~technologies de l'information et de la communication~» fondée sur la classification internationale des activités économiques. La définition retient trois verbes~: \emph{capture, transmit and display data and information electronically}. Elle délimite un périmètre statistique~; elle ne formule pas de fonctions. Le stockage n'y figure pas comme catégorie autonome~; le calcul y est subsumé sous le terme générique de \emph{processing}. Cet acte de nomination a ses propres conditions de possibilité, et leur généalogie éclaire les angles morts que la thèse cherche à combler.

Le mot «~télécommunication~» lui-même n'apparaît qu'en 1904, forgé par Édouard Estaunié, polytechnicien et ingénieur aux Postes et Télégraphes, dans son \emph{Traité pratique de télécommunication électrique}\footnote{\textcite{estaunie_traite_1904}. Estaunié, qui fut aussi romancier et membre de l'Académie française, construit le néologisme à partir du grec \emph{τηλε} (loin) et du latin \emph{communicare} (partager), assumant le mélange des racines au profit de l'intelligibilité. Le geste est de voir que la télégraphie, la téléphonie et les transmissions radio, bien que techniquement hétérogènes, remplissent une même fonction~: communiquer à distance par des moyens électriques.}. Le mot met vingt-huit ans à devenir une catégorie institutionnelle internationale~: c'est la Conférence de Madrid de 1932 qui l'adopte en fusionnant l'Union Télégraphique Internationale (1865) et l'Union Radiotélégraphique (1906) en une Union Internationale des Télécommunications. La définition retenue~--- «~toute transmission, émission ou réception de signes, signaux, écrits, images et sons, ou renseignement de toute nature, par fil, radio, optique ou autres systèmes électromagnétiques~»~--- est une définition de la fonction \emph{transmettre}. Elle ne couvre ni le calcul ni le stockage.

En 1948, Shannon publie \emph{A Mathematical Theory of Communication} et déplace l'objet~: ce n'est plus la transmission qui est définie, c'est le contenu transmis. L'information est quantifiée en bits, déclarée indépendante du sens sémantique et du support physique. L'entropie informationnelle est formellement analogue à l'entropie thermodynamique~--- Shannon reprend le mot à dessein~--- mais les deux grandeurs ne sont pas identifiées~: le bit n'a pas de joule. Le geste de Shannon est le point d'origine de l'une des trois séparations que la présente thèse examine~: la séparation entre information et énergie, qui n'est pas un fait technique mais un choix de formalisation dont les conséquences se propagent dans toute la modélisation économique ultérieure.

Ce que Shannon formalise en 1948 dans le domaine de la communication, Wiener l'avait posé la même année comme principe général de la cybernétique~: tous les phénomènes du monde peuvent être caractérisés comme masse, comme énergie, ou comme information\footnote{\textcite{wiener_cybernetics_1948}, p.~155~: «~Information is information, not matter or energy. No materialism which does not admit this can survive at the present day.~» Ce geste fonde la séparation entre information et énergie que Shannon amplifiera. C'est aussi contre ce geste que Ropohl écrira, trente ans plus tard~: «~en réalité, l'information est fondamentalement liée à des signaux matériels ou énergétiques~» (\textcite{ropohl_systemtheorie_1979}, §~3.2.3).}. C'est \textcite{ropohl_systemtheorie_1979} qui transpose cette tripartition dans une théorie générale des systèmes sociotechniques~: dans tout système d'action, les inputs, états et outputs se caractérisent comme masse, énergie ou information\footnote{\textcite{ropohl_systemtheorie_1979}~; voir aussi \textcite{ropohl_philosophy_1999}. Ropohl, ingénieur et philosophe, tient la tripartition pour un schéma heuristique, pas une propriété ontologique~: en réalité masse et énergie sont équivalentes, et l'information est liée à des signaux physiques. De ce même cadre, il tire une classification des systèmes techniques par deux dimensions~: catégorie d'output (masse, énergie, information) et classe de fonction (transformation, transport, \emph{stockage}). Le stockage y figure comme catégorie analytique autonome~--- confirmé par la liste des sous-fonctions de tout système d'action que \textcite{miller_living_1978} établit (réception, \emph{stockage}, traitement), et validant rétrospectivement l'inclusion de \textsc{store} dans la tripartition fonctionnelle de cette thèse.}. De ce principe, \textcite{hilty_ict4s_2012} tire une périodisation des métaparadigmes techniques~: la transformation de la matière (des âges des matériaux à la métallurgie), la transformation de l'énergie (les révolutions industrielles, de la vapeur à l'électricité), et la transformation de l'information~--- le troisième paradigme, dont nous sommes contemporains\footnote{Que chaque palier de cette périodisation bouleverse la société dans ses fondements, Marx l'avait formulé en 1856 à propos de la transition vers le deuxième~: \og~\emph{Der Dampf, die Elektrizität und der Selfaktor waren revolutionärer als selbst die Bürger Barb\`es, Raspail und Blanqui}~\fg~---~\og~La vapeur, l'électricité et la machine à filer ont été des révolutionnaires plus dangereux que Barb\`es, Raspail et Blanqui eux-mêmes~\fg~(\textcite{marx_speech_1856}, cité par \textcite{ropohl_philosophy_1999}). Ropohl reprend cette intuition pour en faire le principe général de sa philosophie sociotechnique~: la technique n'est pas un destin autonome ni une somme d'artefacts individuels, mais la transformation organisée de la matière, de l'énergie ou de l'information dans des systèmes d'action humains.}. C'est sur cette base qu'il construit son analyse de la puissance transformatrice de l'ICT~: cette puissance vient précisément de la confluence, sur un substrat unique, des trois lignes technologiques qui traitent l'information~--- la transmettre, la stocker, la traiter. Séparées pendant tout le \textsc{xix}\ieme{} siècle et la première moitié du \textsc{xx}\ieme{}, ces trois lignes convergent avec la numérisation~: un signal numérique peut être copié sans perte, stocké, et traité algorithmiquement sur le même substrat. \textcite{hilbert_worlds_2011} datent empiriquement ce basculement à 2002, année où pour la première fois la capacité mondiale de stockage numérique dépasse la capacité analogique~; dès 2012, 99\,\% de l'information stockée dans le monde est numérique.

Chaque acte de nomination unifie et occulte. «~Télécommunication~» (1904) unifie télégraphe et téléphone mais occulte le calcul. «~Information~» (1948) unifie signal et message mais occulte l'énergie. «~ICT~» (1998) unifie télécommunications et informatique mais occulte le stockage comme fonction distincte. La convergence que Ropohl théorise comme troisième métaparadigme et que Hilty célèbre comme source de puissance transformatrice a un revers que ni l'un ni l'autre ne formule~: c'est précisément parce que les trois fonctions convergent sur un substrat unique que le couplage entre information et énergie devient structurel et dominant dans le régime actuel. Avant la convergence numérique, une information stockée sur papier ou sur carton perforé ne consommait de l'énergie qu'à la fabrication et à la lecture~; une information stockée sur substrat numérique consomme de l'énergie en permanence, et cette consommation croît avec le stock. La tripartition fonctionnelle~--- transmettre, stocker, calculer~--- que cette thèse adopte reprend celle de Hilty et Hilbert, mais en renversant le signe~: là où ils voient dans la convergence une augmentation de capacité, cette thèse voit dans la même convergence la genèse d'un couplage énergétique que les cadres analytiques existants ne capturent pas.

Trois traditions proposent des cadres pour penser les transformations techniques qui restructurent les économies, et aucune ne suffit à rendre compte du phénomène que la présente thèse étudie.

La première, née dans les années~1990 autour de \textcite{bresnahan_general_1995} puis systématisée par \textcite{helpman_general_1998} et \textcite{lipsey_economic_2005}, identifie un petit nombre de «~technologies à usage général~» (GPT) dont les propriétés communes~--- pervasivité, dynamisme technologique propre, complémentarités innovantes~--- suffiraient à rendre compte des transformations économiques qu'elles engendrent. La vapeur, l'électricité, l'ordinateur seraient des chocs discrets, identifiables, dont les effets se déploient selon des phases récurrentes. \textcite{lipsey_economic_2005}, dans \emph{Economic Transformations}, poussent le geste le plus loin~: vingt GPT identifiées dans toute l'histoire, une décomposition en cinq catégories structurelles (technologie, structure facilitante, politique, structure politique, performance), trois niveaux d'abstraction successifs conduisant de l'observation historique à la simulation formelle. Le cadre a une vertu~: il nomme les chocs, les périodise, et produit des prédictions~--- un pattern en cinq phases, des courbes logistiques d'efficience et d'application. Mais il a un angle mort structurel~: il traite chaque GPT comme un bloc. «~L'ICT~» est une GPT~; la distinction entre ce que l'ICT transmet, ce qu'elle stocke et ce qu'elle calcule n'a pas de place dans l'architecture du cadre. Or l'examen historique, même sommaire, montre que deux technologies de calcul de la même époque, vendues aux mêmes clients, fondées sur des mécanismes comparables~--- la calculatrice mécanique et la tabulatrice à cartes perforées~--- produisent des structures de marché opposées~: un marché fragmenté d'une centaine de firmes pour la première, un monopole à 85,7\,\% pour la seconde. Si l'on reste au niveau de la GPT comme bloc, cette variance est invisible. La GPT discrète voit les chocs~; elle ne voit pas la structure interne qui détermine leurs effets.

\textcite{vanderkooij_invention_2015}, dans une série de monographies consacrées aux «~engines~» de la GPT-Électricité (machine à vapeur, moteur électrique, télégraphe, téléphone, radio), tente de combler le trou que Lipsey lui-même avait identifié~: les microfondations technologiques de la GPT. Sa réponse est le «~General Purpose Engine~» (GPE), l'artefact technique dont les propriétés (pervasivité, dynamisme, essaimage) sont celles que la littérature attribue à la GPT dans son ensemble. Le GPE est le noyau~; la GPT est la constellation d'innovations qui se forme autour de lui. Le travail de Van der Kooij est remarquable par sa granularité~: il descend au niveau du brevet, de l'ingénieur, du prototype, là où Lipsey reste au niveau de la méta-technologie. Ses monographies fournissent un matériau empirique irremplaçable. Mais son cadre analytique pose deux problèmes. Le premier est la distinction qu'il opère entre «~communication technology~» (CT) et «~information technology~» (IT), qu'il traite comme deux trajectoires parallèles au sein de la GPT-Électricité. Cette partition binaire occulte une troisième fonction~: le stockage. La carte perforée de Hollerith, la mémoire DRAM, le disque dur ne sont ni de la communication ni du traitement au sens strict~; ce sont des supports de persistance, et c'est précisément la persistance qui rend le couplage entre système informationnel et infrastructure physique structurel~--- une information stockée numériquement consomme de l'énergie en permanence, ce qui n'est le cas ni du papier, ni du carton, ni du vinyle. Le second problème est que Van der Kooij ne déploie pas son cadre analytique complet sur la période du calcul mécanique~: Hollerith reçoit une demi-phrase là où Morse, Bell et Marconi reçoivent chacun un volume entier. Ce n'est pas un oubli~; c'est le signe que le \emph{compute}, trop petit relativement à l'économie avant la seconde moitié du vingtième siècle, échappe au radar d'un cadre construit pour les «~engines~» qui transforment visiblement la société. La thèse utilise Van der Kooij jusqu'au point où son cadre casse, et c'est ce point de cassure~--- l'absence du \emph{store}, la marginalisation du \emph{compute}~--- qui définit l'apport propre.

La deuxième tradition est celle du «~development block~», forgé par Dahmen dans les années~1950 et repris par \textcite{nuvolari_understanding_2020} pour comparer les révolutions industrielles successives. L'idée est celle de l'autocatalyse~: un ensemble de composants complémentaires où l'amélioration de chacun stimule l'innovation dans les autres. La vapeur, le fer, le charbon et les machines-outils forment un circuit~; les semi-conducteurs, les ordinateurs, le logiciel et les réseaux en forment un autre. Le concept a une force que la GPT discrète n'a pas~: il voit les complémentarités internes au système technique et il explique pourquoi certaines combinaisons produisent des transformations autocatalytiques et d'autres non. Mais la décomposition de Nuvolari opère par composant technologique, pas par fonction. Dire que le \emph{development block} ICT est fait de semi-conducteurs, d'ordinateurs, de logiciel et de réseaux, c'est nommer les pièces~; ce n'est pas dire ce que chaque pièce fait, ni pourquoi des pièces qui font la même chose~--- calculer~--- par des modes d'appropriation différents produisent des effets économiques opposés.

La troisième tradition mesure plutôt qu'elle ne nomme. \textcite{grubler_rise_1990,grubler_technology_1998} montrent que les clusters transport-communication se diffusent et se saturent ensemble selon des courbes logistiques~; \textcite{fouquet_digitalisation_2023} construit des séries longues d'intensité énergétique et d'intensité informationnelle depuis~1850 et en tire un isoquant énergie-information. Ces travaux ont une vertu décisive~: ils produisent des grandeurs, pas des catégories. Mais ils ne produisent pas de mécanismes. Grübler peut montrer que canaux et alimentation animale saturent ensemble vers~1870, que rail et charbon saturent ensemble vers~1930~; il ne peut pas dire par quel canal économique la saturation de l'un entraîne celle de l'autre. Fouquet peut montrer que le taux marginal de substitution de l'information à l'énergie décroît dans le temps~; il ne peut pas décomposer cet isoquant par fonction informationnelle pour localiser d'où vient la décroissance.

Aucun de ces cadres ne pose la question que la présente thèse constitue en objet~: quel est le rapport entre les fonctions informationnelles~--- transmettre, stocker, calculer~--- et les structures de l'économie, du capital au marché en passant par le travail~? Et quel est le coût de ce rapport en énergie et en matière~? La GPT discrète voit les chocs mais pas la structure fonctionnelle interne qui en détermine les effets. Le \emph{development block} voit les complémentarités mais les décompose par composant technologique, pas par fonction. La transformation continue voit les tendances mais ne formule pas les mécanismes économiques qui les produisent. La tripartition fonctionnelle~--- transmettre, stocker, calculer~--- que \textcite{hilbert_worlds_2011} ont quantifiée comme trois grandeurs mondiales indépendantes, est l'unité d'analyse qui permet de voir simultanément la continuité (les trois fonctions persistent à travers les régimes techniques), la variance (chaque fonction produit des effets économiques distincts) et la formation progressive d'un \emph{development block} fonctionnel dont les composants, pour la première fois dans l'histoire de l'ICT, sont devenus interopérables et autocatalytiques.


\section*{Plan de la thèse}

Le couplage entre ICT et énergie n'appartient à aucune discipline constituée. L'historien des techniques en voit les manifestations~--- les trajectoires d'innovation, les ruptures de régime, les dépendances au sentier~--- mais il ne mesure ni la productivité agrégée ni les flux physiques. Le macroéconomiste observe des anomalies dans la structure de l'investissement, dans les formes de financement, dans la dynamique sectorielle, mais il ne voit pas les mécanismes techniques qui les produisent ni l'empreinte matérielle qui en résulte. L'économiste de l'énergie, enfin, mesure des consommations, identifie des effets rebond, quantifie des trajectoires d'émissions~--- mais les formes institutionnelles qui organisent ces flux lui échappent pour l'essentiel. Chacune de ces perspectives a son domaine de certitude et ses angles morts~; chacune constate des faits que les deux autres ne voient pas, et suppose des liens qu'elle ne peut pas établir seule. Les trois premiers chapitres de cette thèse tirent parti de cette asymétrie plutôt que de la contourner. Chacun se place délibérément dans l'un de ces cadres disciplinaires~--- l'histoire des techniques, la macroéconomie, l'économie de l'énergie~--- et observe de là les deux autres. Le premier chapitre, ancré dans l'ICT, constate des trajectoires techniques et laisse émerger du récit lui-même les questions macroéconomiques et énergétiques auxquelles il ne saurait répondre. Le deuxième, ancré dans l'analyse économique, constate des transformations de structure et pose des questions sur leur empreinte physique que seul le troisième pourra prendre en charge. Le troisième, ancré dans l'énergie, boucle vers les deux premiers en montrant que les contraintes matérielles qu'il documente ne se comprennent qu'à la lumière des mécanismes techniques et économiques établis dans les chapitres précédents. La progression n'est donc pas linéaire~: elle est spiralaire, chaque chapitre enrichissant rétrospectivement la lecture des précédents. Au terme de ce parcours, le couplage ICT-énergie apparaît comme un objet qu'aucune des trois disciplines ne saisit seule, mais que leur croisement rend visible. C'est cette visibilité acquise qui légitime~--- et qui appelle~--- le passage à la modélisation.

\newpage

Le \autoref{ch:history} (\emph{Transmettre, stocker, calculer}) examine, cas par cas et dans l'ordre chronologique, les effets des innovations informationnelles sur les structures économiques et identifie les basculements de régime technique qui les ont portés. Le parcours va des tablettes d'argile sumériennes et du calcul organisé de Prony à l'actuelle convergence d'une infrastructure numérique intégrée~--- télégraphe, chemin de fer, calcul mécanisé, téléphone, radiodiffusion, informatique, Internet, intelligence artificielle~---, en observant chaque cas au prisme de l'économie et du couple énergie-matière. La méthode est inductive~: les questions analytiques émergent du récit lui-même, par abduction au sein de chaque cas et confrontation d'un cas au suivant. Le chapitre se conclut par une réduction typologique qui ordonne les récurrences observées en dimensions économiques et en forces transversales (énergie-matière, souveraineté), et formule l'hypothèse de continuité de la fonction informationnelle à travers les changements de substrat technique.

Le \autoref{ch:economics} (\emph{Analyse économique des structures transformées par l'ICT}) reprend le matériau historique sous un angle économique systématique. Pour chacune des dix dimensions identifiées en fin du chapitre~1~--- propriétés du bien informationnel, rendements d'échelle, nature du capital, travail et division cognitive, demande et préférences, signal de prix et structure des marchés, coûts de transaction et formes organisationnelles, financement et cycles d'infrastructure, géographie et commerce, trajectoires d'innovation et dépendance de sentier~---, le chapitre délimite le champ théorique, identifie les mécanismes causaux, les confronte au matériau historique et aux données contemporaines, et en formule les résultats comme faits stylisés. Chaque dimension engage sa propre tradition théorique~; le chapitre n'est pas un état de l'art mais une analyse qui produit des mécanismes structurels suffisamment précis pour être testables dans un modèle. La section finale examine les interactions entre dimensions et identifie les dynamiques macroéconomiques qui en résultent~: investissement et allocation de l'épargne, distribution fonctionnelle des revenus et demande effective, changement structurel et trajectoire de long terme.

Le \autoref{ch:energy} (\emph{Énergie, climat, ressources}) adopte le point de vue du substrat biophysique. Nourri des observations historiques et des mécanismes économiques des deux chapitres précédents, il confronte la relation entre information et énergie aux données de longue durée~--- séries d'intensité énergétique et informationnelle, co-évolutions des secteurs ICT-énergie, isoquants énergie-information~--- et aux données contemporaines sur la consommation des infrastructures numériques. Le chapitre décompose la relation par fonction (transmission, stockage, calcul) pour montrer que le changement de régime énergétique est localisé dans la couche du calcul intensif en intelligence artificielle, tandis que les couches réseau et stockage classique restent en régime d'efficacité croissante. Il formule l'hypothèse d'un changement de régime dans le couplage ICT-énergie et en examine les conditions, la portée et les limites.

Le \autoref{ch:imaclim} (\emph{Dispositif expérimental}) construit le cadre de modélisation~: analyse critique des stratégies de représentation du numérique dans les modèles d'évaluation intégrée, justification du cadre IMACLIM~--- modèle hybride articulant représentation macroéconomique en valeur et représentation physique des flux d'énergie, adapté aux situations de second-best~--- puis construction des scénarios et de leurs hypothèses structurantes. Trois variables sont endogénéisées~: le capital computationnel-énergétique, la puissance électrique contractée par les opérateurs de centres de données, et l'intensité énergétique du calcul.

Le \autoref{ch:scenarios} présente les résultats et les interprète à la lumière des trois premiers chapitres.

Le \autoref{ch:results} discute les implications.

Le \autoref{ch:discussion} ouvre des perspectives.